\providecommand{\Lie}{\mathrm{Lie}}
\section{Factorisation-algebra gluing}
\label{sec:gluing:factorisation}
\label{part:gluing-factorisation}
The gluing modes developed in Parts~II--IV record, chart-wise, the Hall
datum of a local CoHA and transmit it across overlaps by a cocycle whose
values are Fourier--Mukai kernels, Bridgeland--King--Reid equivalences,
or Van den Bergh tilting modules. The gluing mode developed now records
the chart-wise datum as a \emph{factorisation algebra} in the sense of
Costello--Gwilliam \cite{CostelloGwilliam,CostelloGwilliam2017} and
transmits it by the Weiss cosheaf axiom of \cite{CostelloGwilliam}
Vol.~I Definitions~$6.1.0.5$ \textup{(}Weiss cover\textup{),}~$6.1.0.6$
\textup{(}Weiss refinement\textup{),} and~$6.1.2.1$ \textup{(}homotopy-
cosheaf/factorisation axiom\textup{).} The two presentations will
coincide, for toric CY$_3$ and for K3 $\times$ E, at the level of
sections of the assembled sheaf; the factorisation-algebra framing
exposes the piece of the data the other modes silence, namely the
\emph{global} compatibility along the Ran space of the reference curve.

Four sections. Section~\ref{sec:fa-weiss-covers} defines a Weiss cover,
derives why a two-chart open cover is never a Weiss cover, and records
the subordinate polydisc refinement that turns any open cover into a
Weiss cover of the form appropriate to factorisation-algebra locality.
Section~\ref{sec:fa-ran-bd} constructs the Ran space $\Ran(C)$ of a
smooth complex curve $C$ as a colimit of symmetric products and names
the Beilinson--Drinfeld chiral algebra as a factorisation cosheaf on
$\Ran(C)$, together with the Francis--Gaitsgory
\cite{FrancisGaitsgory2012} chiral Koszul duality and its
higher-dimensional $E_n$-analogue \cite{Francis2013}; it also isolates
finite DWR/Ran Rees comparison as an ordered $E_1$ Hall descent, with
the $E_2$ centre or double left to the separate Drinfeld-centre
operation.
Section~\ref{sec:fa-two-stage} derives the two-stage factorisation
$\Phi_3 = \SpCh_{\Sigma_2, C} \circ \PhiFA_3$ from first principles,
reading Stage~$1$ as a factorisation-algebra assembly on the CY$_3$
target and Stage~$2$ as a specialisation by factorisation homology
over a transverse $E_2$-fibre. Section~\ref{sec:fa-hol-top} separates
the holomorphic and topological factors of 5D holomorphic Chern--Simons
on $\R \times \C^2$ and names the stratified manifold structure
(Costello--Gwilliam Volume~II~\cite{CostelloGwilliam2017}~\S$4.10$) along
which the two factor-wise factorisation algebras glue.

The super-factorisation axiom on a super-factorisation algebra valued
in $\Ch(\sVect)$ is a \emph{strict} equality of structure maps after
Koszul super-interchange, not a coherence datum: the target
$\Ch(\sVect)$ is $1$-categorical symmetric monoidal, with associator and
symmetry equalities of functors rather than higher coherences. The
$\Z/2$-grading on the super-Yangian arising from the conifold two-chart
gluing is therefore read off an equality of factorisation sections, not
a homotopy datum. This $1$-categorical super-interchange and the
$(\infty, 1)$-categorical Ran-local refinement of factorisation locality
belong to different layers of the construction.

\subsection{Weiss covers}
\label{sec:fa-weiss-covers}

\subsubsection{The factorisation local-to-global axiom}

Fix $M$ a smooth manifold (the eventual use will take $M$ to be either
a smooth complex curve $C$, a smooth complex surface, or a smooth
complex threefold; the definition is uniform across these cases).
A \emph{prefactorisation algebra} $\cF$ on $M$ valued in chain
complexes over $k$ \cite{CostelloGwilliam} Definition~3.1.1 assigns to
each open $U \subset M$ a chain complex $\cF(U)$ and to each finite
collection of pairwise-disjoint opens $U_1, \dots, U_n \subset V$
inside an open $V \subset M$ a structure map
\begin{equation}
\label{eq:prefactorisation-structure-map}
 m_{U_1, \dots, U_n; V} \colon
 \cF(U_1) \otimes \cdots \otimes \cF(U_n) \longrightarrow \cF(V)
\end{equation}
subject to the associativity and symmetry axioms of a prefactorisation
algebra. A prefactorisation algebra is a \emph{factorisation algebra}
if, in addition, it satisfies the \emph{factorisation local-to-global
axiom}: for every open $V \subset M$ and every \emph{Weiss cover}
$\fU$ of $V$, the natural map from the homotopy colimit of the
Čech complex of $\cF$ over $\fU$ to $\cF(V)$ is an equivalence.

The axiom is useless without a definition of Weiss cover that makes it
hold for the examples of interest; the combinatorics of Weiss covers
is what distinguishes factorisation-algebra locality from ordinary
sheaf locality.

\begin{definition}[Weiss cover; \protect{\cite{CostelloGwilliam}}
 Vol.~I Definition~$6.1.0.5$]
\label{gl:def:weiss-cover}
Let $V \subset M$ be an open subset. A collection $\fU = \{U_\alpha\}$
of open subsets $U_\alpha \subset V$ is a \emph{Weiss cover} of $V$ if,
for every finite set of points $\{x_1, \ldots, x_k\} \subset V$, there
exists an index $\alpha$ with $\{x_1, \ldots, x_k\} \subset U_\alpha$.
The Weiss covers of $V$ generate the \emph{Weiss Grothendieck
topology} on $\mathrm{Opens}(V)$.
\end{definition}

Equivalently: the cover $\fU$ is closed under finite disjoint unions
inside $V$, up to refinement by a cofinal subset. The formulation in
terms of finite subsets is operational: to check that $\fU$ is a Weiss
cover one exhibits, for each finite point-configuration $S$ in $V$,
one open $U_\alpha \in \fU$ containing $S$.

The reason the axiom demands Weiss covers and not ordinary open covers
is the following. A factorisation algebra $\cF$ records the \emph{joint}
algebraic data of observables localised near finite point-configurations.
Two observables supported near disjoint point-configurations $S_1 \sqcup
S_2$ must multiply through the structure map $m_{U_1, U_2; V}$ where
$U_1 \supset S_1$ and $U_2 \supset S_2$ are disjoint enlargements; the
structure is reconstructed locally only if the cover contains opens
separating every such pair. Ordinary open covers need not; Weiss covers
do, by the finite-point-configuration criterion.

\subsubsection{Why a two-chart cover is not a Weiss cover}
\label{subsec:two-chart-not-weiss}

The sharpest elementary illustration arises in the conifold two-chart
atlas of \S\ref{sec:gluing:toric-cech}. Let $Y = U_+ \cup U_-$ with $U_\pm \cong \C^3$ and
$U_+ \cap U_- \cong \G_m \times \C^2$. One asks whether $\fU_0 = \{U_+,
U_-\}$ is a Weiss cover of $Y$.

Take a two-point configuration $S = \{p_+, p_-\} \subset Y$ with
$p_+ \in U_+ \setminus U_-$ and $p_- \in U_- \setminus U_+$; such
configurations exist because $U_+$ and $U_-$ are each proper subsets
of $Y$ (neither chart covers the whole variety). The finite subset
$S$ is contained in neither $U_+$ nor $U_-$: $p_- \notin U_+$ rules out
$U_+$, and $p_+ \notin U_-$ rules out $U_-$. Consequently no member of
$\fU_0$ contains $S$, and $\fU_0$ fails the finite-subset criterion of
Definition~\ref{gl:def:weiss-cover}.

\begin{proposition}[Two-chart atlases do not satisfy Weiss locality]
\label{prop:two-chart-not-weiss}
Let $M$ be a connected smooth manifold of positive dimension and
$\fU = \{U_1, U_2\}$ a two-element open cover of $M$ with
$U_1 \setminus U_2 \neq \varnothing$ and $U_2 \setminus U_1 \neq
\varnothing$. Then $\fU$ is not a Weiss cover of $M$.
\end{proposition}

\begin{proof}
Choose $p_1 \in U_1 \setminus U_2$ and $p_2 \in U_2 \setminus U_1$.
The two-point configuration $S = \{p_1, p_2\}$ is not contained in
either $U_1$ (because $p_2 \notin U_1$) or $U_2$ (because $p_1 \notin
U_2$), nor in any other member of $\fU$ (there are none). The
finite-subset criterion fails at $S$.
\end{proof}

The conifold two-chart cover is the prototype: every Čech-style
two-chart gluing in \S\ref{sec:gluing:toric-cech} produces a two-element cover failing Weiss
locality. Gluing by factorisation-algebra locality therefore requires
a refinement of the two-chart cover to a Weiss cover, and the
refinement is what injects the \emph{finitely many} many-point
configurations which the Čech nerve of the two-chart atlas cannot see.

\subsubsection{Weiss refinement procedure}
\label{subsec:weiss-refinement}

The refinement procedure is explicit. Given an open cover $\fU_0 =
\{U_\alpha\}$ of $M$, choose the subordinate polydisc basis
\[
 \mathfrak B(\fU_0)
 :=
 \{\,D\subset M \mid D \text{ is a sufficiently small coordinate
 polydisc and } \overline D\subset U_\alpha
 \text{ for some } \alpha\,\}.
\]
Form the \emph{disjoint-union closure}
\begin{equation}
\label{eq:weiss-refinement}
 \mathfrak B(\fU_0)^{\sqcup} \;:=\;
 \bigl\{\,
 D_1 \sqcup \cdots \sqcup D_n
 \;\bigm|\; n \ge 1,\; D_i \in \mathfrak B(\fU_0),\;
 D_i \cap D_j = \varnothing \text{ for } i \ne j
 \,\bigr\}.
\end{equation}
Every finite subset $S = \{p_1, \dots, p_m\} \subset M$ admits, by
first choosing $U_{\alpha_i}\ni p_i$ and then shrinking inside
coordinate charts, pairwise-disjoint polydiscs
$D_i\in\mathfrak B(\fU_0)$ with $p_i\in D_i$. The disjoint union
$D_1\sqcup\cdots\sqcup D_m\in\mathfrak B(\fU_0)^{\sqcup}$ contains
$S$; hence $\mathfrak B(\fU_0)^{\sqcup}$ is a Weiss cover. The
subordinate-polydisc step is essential: the disjoint-union closure of
the original members of $\fU_0$ need not add any open set.

\begin{remark}[Weiss refinement of the conifold two-chart atlas]
\label{rem:weiss-refinement-conifold}
For the conifold $Y = U_+ \cup U_-$, the disjoint-union closure of the
original two chart opens reduces to $\fU_0$ itself:
the only disjoint families of members of $\fU_0$ are the two singletons
$\{U_+\}$ and $\{U_-\}$ (the pair $U_+, U_-$ is not disjoint since
$U_+ \cap U_- \cong \G_m \times \C^2 \ne \varnothing$), and the
disjoint-union of a singleton equals its single member. The
disjoint-union closure produces no new opens, and
Proposition~\ref{prop:two-chart-not-weiss} continues to apply to
$\fU_0$. The refinement therefore requires enlarging
$\fU_0$ with subordinate opens sharper than the two charts: for each
$p \in Y$ choose
a contractible open neighbourhood $D_p \subset U_\pm$ (whichever chart
contains $p$), small enough that finitely many such $D_p$ can be made
pairwise disjoint by shrinkage. The family $\fD = \{D_p : p \in Y\}$
is an ordinary open cover of $Y$; its disjoint-union closure $\fD^{\sqcup}$
is a Weiss cover of $Y$ by the argument following
equation~\eqref{eq:weiss-refinement}. The original two-chart cocycle
is pulled back to the \v{C}ech nerve of $\fD^{\sqcup}$, which keeps
track of all finite multi-point configurations, including the pairs in
$U_+ \cap U_-$.
\end{remark}

The refinement produces, in particular, Weiss-cover cells indexed by
finite point-configurations rather than pairs of charts. Its ordered
lift is indexed by $\Conf_n(Y)$, $n \geq 1$, while its unordered
quotient maps to the Ran stratification; this is the
factorisation-algebra analogue of the Ran-space decomposition of
Section~\ref{sec:fa-ran-bd}. The ordered lift is the layer on which
Hall multiplication remembers the direction of short exact sequences.

\subsection{Ran space and Beilinson--Drinfeld chiral algebras}
\label{sec:fa-ran-bd}

\subsubsection{The Ran space of a curve}

Let $C$ be a smooth complex curve. The \emph{Ran space} of $C$ is the
set of non-empty finite subsets of $C$ topologised by the colimit
\begin{equation}
\label{eq:ran-space-colimit}
 \Ran(C) \;=\; \varinjlim_{\, n \ge 1\,} \; C^n / S_n
\end{equation}
in the category of (ind-)topological spaces, where the transition maps
$C^n/S_n \to C^{n+1}/S_{n+1}$ send a tuple $(x_1, \dots, x_n)$ to
$(x_1, x_1, x_2, \dots, x_n)$ (repeating the first entry). The
presentation \eqref{eq:ran-space-colimit} identifies configurations
differing only by repetition, so $\Ran(C)$ parameterises finite subsets
of $C$ regardless of multiplicity.

The Ran space is the natural index of factorisation data on $C$. A
prefactorisation algebra on $C$ assigns observables to finite
point-configurations; the assignment extends along the transition maps
\eqref{eq:ran-space-colimit} by the structure maps
\eqref{eq:prefactorisation-structure-map} (at repeated points, the
structure map at a pair of coinciding points specialises to the
diagonal). The factorisation local-to-global axiom, specialised to
$M = C$, becomes the condition that the output factorisation cosheaf
on $\Ran(C)$ satisfies the homotopy-cosheaf axiom
\cite{CostelloGwilliam} Vol.~I Definition~$6.1.2.1$ evaluated on the
Ran site generated by Weiss covers of the symmetric-product stratum
$C^n/S_n$; equivalently, by Francis--Gaitsgory
\cite{FrancisGaitsgory2012} Theorem~$5.2.1$, the cosheaf is supported on
the diagonal stratification of $\Ran(C)$, with factorisation data
determined by restriction to each multi-point stratum.

\subsubsection{Finite DWR/Ran Rees comparison}
\label{subsec:finite-dwr-ran-rees-comparison}

The Weiss condition is infinite in arity, but every perturbative
hCS--Hall comparison is tested at finite height before any completed
limit is taken. Fix a subordinate DWR polydisc refinement
$\mathfrak U$, a charge monoid $\Gamma$, and finite bounds
\[
 N \quad\text{(holomorphic jets)},\qquad
 r \quad\text{(renormalisation order)},\qquad
 L \quad\text{(charge)},\qquad
 m \quad\text{(Ran arity)}.
\]
Let
\[
 \mathsf N^{\Ran,\ord}_{\le m}(\mathfrak U;\Gamma_{\le L})
\]
be a chosen finite full, refinement-closed sub-nerve whose $p$-simplices
are ordered finite families
\[
 \sigma=(s_1<\cdots<s_q;\,P_{s_i}\Subset
 U_{i,0}\cap\cdots\cap U_{i,p},\,\gamma_{s_i}\in\Gamma_{\le L}),
 \qquad q\le m,
\]
with the $P_{s_i}$ pairwise disjoint subordinate polydiscs. Faces and
refinements are generated by deleting an overlap, shrinking a
polydisc, and passing to a subordinate intersection; refinement-closed
means that every higher simplex has a vertex refinement inside the same
finite sub-nerve. This is the ordered lift of the finite Ran stratum;
the order is not a time parameter, but the Hall order on short exact
sequences.

\begin{proposition}[Finite DWR/Ran Rees descent and the ordered Hall layer]
\label{prop:finite-dwr-ran-rees-ordered-hall-layer}
For every simplex $\sigma\in
\mathsf N^{\Ran,\ord}_{\le m}(\mathfrak U;\Gamma_{\le L})$, assume
the following data.
\begin{enumerate}
\item A finite hCS complex
$\Obs_{\mathrm{hCS}}^{q,\le r,\le N}(\sigma)_{\le L}$ with its
compact-support refinement maps.
\item A finite cyclic critical model
$(V_\sigma,W_\sigma,\omega_\sigma)$ over
$\Omega^\bullet(\Delta^p)$ and the Rees Hall complex
\[
 \Hall_{\sigma}^{\mathrm{Rees},\mathrm{or}}
 =
 \left(
 C^\bullet\!\left(
 \mathfrak g_\sigma,
 \Gamma(V_\sigma,\wedge^\bullet T_{V_\sigma})[[\hbar]]
 \otimes\Omega^\bullet(\Delta^p)
 \right)^{G_\sigma},
 d_{\CE}+\iota_{d_VW_\sigma}
 +\hbar\Delta_{\omega_\sigma}+d_{\Delta^p}
 \right)
 \otimes o_\sigma[s_\sigma](t_\sigma).
\]
\item A continuous degree-zero relative map
\[
 \theta_\sigma^{\mathrm{rel}}\colon
 B\Obs_{\mathrm{hCS}}^{q,\le r,\le N}(\sigma)_{\le L}
 \longrightarrow
 \Hall_{\sigma}^{\mathrm{Rees},\mathrm{or}}
\]
compatible with faces, refinements, orientations, shifts, and Tate
twists.
\end{enumerate}
Assume also multiplicativity on disjoint ordered families:
\[
 \Obs(\sigma\sqcup\tau)\simeq
 \Obs(\sigma)\widehat\otimes\Obs(\tau),\qquad
 W_{\sigma\sqcup\tau}=W_\sigma\boxplus W_\tau,\qquad
 \omega_{\sigma\sqcup\tau}=\omega_\sigma\oplus\omega_\tau,
\]
with $o_{\sigma\sqcup\tau}\simeq o_\sigma\otimes o_\tau$ and
$[s_{\sigma\sqcup\tau}](t_{\sigma\sqcup\tau})=
[s_\sigma+s_\tau](t_\sigma+t_\tau)$. Then
\[
 \Theta_{\mathrm{hCS}\to\Hall}^{\mathrm{Rees},\mathrm{or};N,r,L,m}
 :=
 \sum_{\sigma\in\mathsf N^{\Ran,\ord}_{\le m}(\mathfrak U;\Gamma_{\le L})}
 \int_{\Delta^{|\sigma|}}\theta_\sigma^{\mathrm{rel}}
\]
is a multiplicative natural transformation on the finite DWR/Ran nerve.
If the vertex maps are quasi-isomorphisms, it induces a
quasi-isomorphism on finite DWR/Ran descent.

For disjoint $\sigma$ and $\tau$, the Hall product appearing in this
statement is the Rees Thom--Sebastiani product
\[
 \mathrm{DR}_{W_\sigma,\hbar}\widehat\otimes
 \mathrm{DR}_{W_\tau,\hbar}
 \xrightarrow{\ \mathrm{TS}\ }
 \mathrm{DR}_{W_\sigma\boxplus W_\tau,\hbar},
 \qquad
 \mathrm{DR}_{W,\hbar}:=
 \bigl(\Gamma(V,\wedge^\bullet T_V)[[\hbar]],
 \iota_{dW}+\hbar\Delta_\omega).
\]
Thus the descended object is an ordered $E_1$ Hall factorisation
cosheaf. An $E_2$ object is obtained only after applying the Drinfeld
centre to $\Rep^{E_1}$ of this ordered layer, or after constructing a
completed Drinfeld double from a non-degenerate Hall pairing and
compatible completion. Neither operation is part of the finite
DWR/Ran descent map.
\end{proposition}

\begin{proof}
The relative map $\theta_\sigma^{\mathrm{rel}}$ is a chain map over
$\Omega^\bullet(\Delta^p)$ precisely when it solves the relative
Maurer--Cartan equation in the bar--Rees convolution Lie algebra:
\[
 (d_{\Hall}+d_{\Delta^p})\theta_\sigma^{\mathrm{rel}}
 -\theta_\sigma^{\mathrm{rel}}d_{B\Obs}
 +\frac{1}{2}
 [\theta_\sigma^{\mathrm{rel}},\theta_\sigma^{\mathrm{rel}}]_{\mathrm{conv}}
 =0.
\]
Integrating this identity over $\Delta^p$ and applying Stokes' formula
turns the $d_{\Delta^p}$ term into the alternating sum over faces.
Because the maps are compatible with faces and refinements, these
boundary terms are exactly the \v{C}ech/Ran differential of the finite
sub-nerve. Summing over all simplices gives a Maurer--Cartan element in
the finite total convolution complex, equivalently a natural
transformation on
$\mathsf N^{\Ran,\ord}_{\le m}(\mathfrak U;\Gamma_{\le L})$.

The finite sub-nerve is a finite diagram of complexes. Objectwise
quasi-isomorphisms on vertices extend to objectwise quasi-isomorphisms
on all simplices by the refinement compatibility built into the datum.
Finite totalisations, or equivalently finite homotopy colimits of
complexes, preserve quasi-isomorphisms; the induced descent map is
therefore a quasi-isomorphism.

For multiplicativity, disjointness gives the tensor-product
decomposition of hCS observables: no Feynman graph has an edge joining
two disjoint supports. On the Hall side, the sum potential
$W_\sigma\boxplus W_\tau$ has differential
$\iota_{dW_\sigma}+\hbar\Delta_{\omega_\sigma}$ on the first factor
plus $\iota_{dW_\tau}+\hbar\Delta_{\omega_\tau}$ on the second. Hence
the Rees twisted de Rham complex of the sum is the completed tensor
product of the two Rees twisted de Rham complexes. This is the
Thom--Sebastiani isomorphism for vanishing cycles in Rees form,
compatible with the orientation and Tate factors by the hypotheses
and with the Hall correspondence product by the
Kontsevich--Soibelman construction of critical CoHA
\cite{KS2008}. The product square for $\Theta$ commutes.

The ordered set $\{s_1<\cdots<s_q\}$ records the order of extensions in
the Hall correspondence, so the algebra produced by the proposition is
associative $E_1$ data. No half-braiding, universal $R$-matrix, or
positive--negative Hopf pairing occurs in the hypotheses. Those are
exactly the extra structures supplied by the Drinfeld centre or by a
completed Drinfeld double; hence they cannot be consequences of the
finite DWR/Ran descent alone.
\end{proof}

\subsubsection{Beilinson--Drinfeld chiral algebras as factorisation cosheaves}

A \emph{chiral algebra} on $C$ in the sense of Beilinson--Drinfeld
\cite{BeilinsonDrinfeld} is a $\mathcal{D}$-module $\cA$ on $C$
equipped with a chiral bracket
\begin{equation}
\label{eq:bd-chiral-bracket}
 \mu \colon j_* j^{*} (\cA \boxtimes \cA) \longrightarrow
 \Delta_{!} \cA
\end{equation}
on the complement $j \colon C^2 \setminus \Delta \hookrightarrow C^2$
of the diagonal $\Delta \colon C \hookrightarrow C^2$, subject to a
symmetry axiom and a Jacobi-type associativity axiom
(\cite{BeilinsonDrinfeld}~\S3.4). The bracket $\mu$ quantifies the
short-distance singularity of the operator product expansion of two
chiral operators on $C$ and is the structure one expects in any
conformal-field-theoretic construction on a Riemann surface.

The reformulation in factorisation-algebra terms is the content of
Francis--Gaitsgory \cite{FrancisGaitsgory2012}~\S3.2. To a chiral
algebra $\cA$ on $C$ one associates a factorisation cosheaf
$\cF_\cA$ on $\Ran(C)$ by the \emph{factorisation envelope}
construction: at a finite configuration $\{x_1, \dots, x_n\} \in
\Ran(C)$ the sections of $\cF_\cA$ are given by
\begin{equation}
\label{eq:fact-envelope-sections}
 \cF_\cA(\{x_1, \dots, x_n\}) \;=\; \bigotimes_{i = 1}^{n} \cA|_{x_i},
\end{equation}
the restriction of $\cA$ to each point and the external tensor product
over the configuration. The chiral bracket $\mu$ equips $\cF_\cA$
with the factorisation-cosheaf structure: at a configuration with
coinciding points the bracket \eqref{eq:bd-chiral-bracket} supplies
the specialisation map; at distinct points the structure is the
tensor-product factorisation of \eqref{eq:fact-envelope-sections}.

\begin{theorem}[Francis--Gaitsgory chiral Koszul duality;
 \protect{\cite{FrancisGaitsgory2012}} Theorem~$1.2.4$, with support
 criterion Theorem~$5.2.1$; Beilinson--Drinfeld
 \protect{\cite{BeilinsonDrinfeld}}~\S$3.4.11$ and Theorem~$3.4.9$]
\label{thm:fg-chiral-koszul}
There is an $(\infty, 1)$-categorical Koszul-duality equivalence, on a
smooth complex curve $C$, between chiral Lie algebras on $\Ran C$
supported on $C$ and factorisation coalgebras on $\Ran(C)$
\textup{(}\cite{FrancisGaitsgory2012} Theorem~$1.2.4$\textup{):}
\begin{equation}
\label{eq:fg-koszul-duality-raw}
 \mathrm{Bar}^{\mathrm{ch}} \colon
 \Lie^{\ast}_{\mathrm{nu}}(C) \;\xrightleftharpoons{\simeq}\;
 \coCom^{\mathrm{fact}}_{\mathrm{nu}}(\Ran C) \;\colon\; \Omega^{\mathrm{ch}},
\end{equation}
with factorisation support on $\Ran(C)$ characterised by the diagonal
stratification \textup{(}\cite{FrancisGaitsgory2012} Theorem~$5.2.1$:
a factorisation coalgebra on $\Ran(C)$ is determined by its restriction
to each open $(\Ran C)^{\le n}$ and the compatibility maps across
strata\textup{).} Under the Beilinson--Drinfeld equivalence between
Lie-$\ast$ algebras on $C$ and chiral algebras on $C$
\textup{(}\cite{BeilinsonDrinfeld}~\S$3.4.11$, Theorem~$3.4.9$; chiral
envelope \S$3.7$\textup{),} the factorisation-envelope functor $\cA
\mapsto \cF_\cA$ identifies chiral algebras on $C$ with non-unital
factorisation cosheaves on $\Ran(C)$, and the factorisation homology
of $\cA$ is the Koszul-dual object:
\begin{equation}
\label{eq:fg-koszul-duality}
 \int_{\Ran(C)} \cF_\cA \;\simeq\; C_\bullet^{\mathrm{chiral}}(\cA),
\end{equation}
with the right-hand side the chiral chain complex of
\cite{BeilinsonDrinfeld}~\S4.2.
\end{theorem}

This equivalence repackages the chiral-algebra calculus of
\cite{BeilinsonDrinfeld} as the factorisation-algebra calculus of
\cite{CostelloGwilliam,CostelloGwilliam2017} on a curve. The two
frameworks are equivalent; each makes different operations convenient.
The chiral-algebra framing produces, by factorisation envelope, the
vertex-algebra image one expects in conformal field theory. The
factorisation-algebra framing produces the gluing-cocycle presentation
one expects in the present chapter: chiral observables on $C$ are
reconstructed from their local sections on Weiss covers of $\Ran(C)$.

\subsubsection{Higher-dimensional variants}
\label{subsec:higher-en-fa}

The Francis--Gaitsgory equivalence of Theorem~\ref{thm:fg-chiral-koszul}
admits an $E_n$-analogue due to Francis \cite{Francis2013}. Let $M$
be a smooth manifold of dimension $n$. An \emph{$E_n$-factorisation
algebra} on $M$ is a factorisation algebra whose structure maps
\eqref{eq:prefactorisation-structure-map} are indexed by
$n$-dimensional disc configurations rather than $0$-dimensional point
configurations; equivalently, a factorisation cosheaf on the
$n$-dimensional Ran space
\begin{equation}
\label{eq:higher-ran-space}
 \Ran_n(M) \;=\; \varinjlim_{n_{\mathrm{discs}} \ge 1}
 \Conf_{n_{\mathrm{discs}}}(M) \big/ S_{n_{\mathrm{discs}}}
\end{equation}
with $\Conf_{n_{\mathrm{discs}}}(M)$ the configuration space of
$n_{\mathrm{discs}}$ pairwise-disjoint $n$-discs in $M$.

\begin{theorem}[Francis $E_n$-Koszul duality;
 \protect{\cite{Francis2013}} Theorem~1.1]
\label{thm:francis-en-koszul}
There is an equivalence between the $(\infty, 1)$-category of
$E_n$-factorisation algebras on a smooth $n$-manifold $M$ and the
$(\infty, 1)$-category of $E_n$-algebras in $\Ch(k)$ augmented over
$k$, given by factorisation-envelope integration.
\end{theorem}

This is the tool required in Section~\ref{sec:fa-two-stage}: a CY$_3$
target $X$ is a smooth real $6$-manifold, and the $E_3$-factorisation
algebra $\PhiFA_3(\cC)$ attached to a CY$_3$ category $\cC$ is,
by Theorem~\ref{thm:francis-en-koszul}, equivalently an $E_3$-algebra
in $\Ch(k)$ augmented over $k$. Costello--Li
\cite{CostelloLi2016,CostelloLi2020} refine the $E_3$-structure to its
\emph{holomorphic} incarnation on the complex threefold $X$ via the
BV-quantisation of 6D holomorphic Chern--Simons.

\begin{remark}[Strict super-factorisation in $\Ch(\sVect)$]
\label{rem:strict-super-factorisation}
When the target chain complexes are super-vector-space-valued, the
braiding in the structure map
\eqref{eq:prefactorisation-structure-map} is the Koszul super-braiding
$\sigma_{V, W} \colon V \otimes W \to W \otimes V$,
$v \otimes w \mapsto (-1)^{|v||w|}\, w \otimes v$, computed from the
$\Z/2$-grading of $\sVect = \Z/2\text{-}\Vect_k$. The symmetric
monoidal category $(\Ch(\sVect), \otimes_k, k^{1|0}, \alpha, \sigma)$
is a strict $1$-categorical symmetric monoidal category: the
associator $\alpha_{A, B, C} \colon (A \otimes B) \otimes C
\xrightarrow{\;\sim\;} A \otimes (B \otimes C)$ and the symmetry
$\sigma$ are equalities (not higher homotopies) of chain-complex
morphisms, since $\Ch(\sVect)$ is built component-wise from the
$1$-category $\sVect$ and the Koszul sign rule pins $\sigma$ uniquely
as the canonical natural transformation (Deligne, \emph{Cat\'egories
tensorielles}, Moscow Math.~J.~$2002$, \S$0.1$; Costello--Gwilliam
Vol.~$1$~\S$A$~\cite{CostelloGwilliam}). Consequently the
super-factorisation axiom for a super-factorisation algebra $\cF$
valued in $\Ch(\sVect)$; the compatibility of the structure maps
$m_{U_1, \ldots, U_n; V}$ under the super-interchange $\sigma$ on the
left; is a \emph{strict equality} of chain maps in the
$1$-categorical sense, not a coherence datum parameterised by an
$(\infty, 1)$-mapping space.

The strict super-braiding and the $(\infty, 1)$-categorical
Ran-locality refinement live on different pieces of the construction.
The Francis--Gaitsgory
equivalence of Theorem~\ref{thm:fg-chiral-koszul} and the Francis
$E_n$-equivalence of Theorem~\ref{thm:francis-en-koszul} are
$(\infty, 1)$-categorical in source and target (Lie-algebra and
factorisation-coalgebra $(\infty, 1)$-categories), but the
super-braiding $\sigma$ in the target is $1$-categorical strict,
inherited from $\sVect$; the two levels act on different parts of the
structure and do not collide. The $\Z/2$-grading on the super-Yangian
arising from the conifold gluing (\S\ref{sec:gluing:toric-cech},
\S$2.1$) is read off a strict equality of factorisation sections under
Koszul super-braiding, and does not require higher-categorical
coherence data.
\end{remark}

\subsection{The two-stage factorisation \texorpdfstring{$\Phi_3 = \SpCh_{\Sigma_2, C} \circ \PhiFA_3$}{Phi3 = SpCh o PhiFA3}}
\label{sec:fa-two-stage}

\subsubsection{Stage~1: $E_3$-factorisation algebra on a CY$_3$ target}
\label{subsec:fa-stage-one}

Fix a CY$_3$ category $\cC$ and a smooth complex CY$_3$ variety $X$
with Calabi--Yau form $\Omega_X \in H^0(X, K_X)$. Stage~1 of the
two-stage factorisation assembles an $E_3$-holomorphic factorisation
algebra $\PhiFA_3(\cC)$ on $X$ from the Hochschild-cochain $E_3$-algebra
of $\cC$. The assembly proceeds in three steps, each elementary but
each drawing on a distinct piece of first-principles input.

\paragraph{Step (a): Kontsevich--Tamarkin formality.}
The Hochschild cochain complex $\HH^\bullet(\cC)$ carries a Gerstenhaber
bracket of degree $1 - d$, descending on cohomology from the
composition-circle $\circ$ of cochain operations
(Deligne conjecture; Kontsevich \cite{Kontsevich1997,Kontsevich1999},
Tamarkin $1998$ preprint arXiv:math/$9803025$ / \cite{Tamarkin2007}).
At $d = 3$ the bracket has degree $1 - 3 = -2$. The $E_d$-formality
theorem asserts that under hypotheses $(H1)$--$(H3)$ of
Chapter~\ref{ch:cy-to-chiral}, the Gerstenhaber bracket lifts, up to
the contractible space of choices controlled by the
Grothendieck--Teichm\"uller group $\GRT_1(\Q)$, to an $E_d$-algebra
structure on $\HH^\bullet(\cC)$. The lift is canonical up to the
$\GRT_1(\Q)$-torsor of Drinfeld associators; any two lifts are
connected by an $(\infty, 1)$-isomorphism in the $E_d$-algebra mapping
space. At $d = 3$ the lift supplies an $E_3$-algebra in $\Ch(k)$.
\emph{Super caveat}: in the super-graded case $\HH^\bullet(\cC_{\sVect})$,
Ginzburg--Schedler (arXiv:$0807.0174$) establish super-formality
only at $E_2$, not at $E_3$; the $E_3$-lift on super Hochschild cochains
is open, bounding the applicability of Step~(a) to the bosonic
Gerstenhaber side unless a super-extension is supplied by hand.

\paragraph{Step (b): Costello--Gwilliam factorisation envelope on $X$.}
An $E_3$-algebra $A$ in $\Ch(k)$ assembles, by the factorisation-envelope
construction of \cite{CostelloGwilliam} Vol.~I~\S$3.6$ and Theorem
$3.4.0.1$ (locally constant factorisation algebras on framed
$n$-manifolds are equivalent to $E_n$-algebras;
compare Lurie \cite{LurieHA}~\S$5.5.4$, Francis \cite{Francis2013}
Theorem~$1.1$), into
a topological $E_3$-factorisation algebra $\mathcal{E}_3(A)$ on any
smooth framed real $6$-manifold. The assignment
$A \mapsto \mathcal{E}_3(A)$ is an $(\infty, 1)$-functorial equivalence
between $E_3$-algebras in $\Ch(k)$ and framed locally constant
$E_3$-factorisation algebras on $\R^6$, and extends by Weiss
local-to-global to an $E_3$-factorisation algebra on any framed smooth
real $6$-manifold $M \cong X_{\mathrm{top}}$. At this step the algebra
is still topological: it uses only the framed real $6$-manifold
structure of $X$, ignoring the complex structure on $X$.

\paragraph{Step (c): Costello--Li holomorphic refinement via $\Omega_X$.}
The Calabi--Yau form $\Omega_X \in H^0(X, K_X)$ supplies, by contraction
with the Dolbeault operator $\bar\partial_X$, the BV pairing of 6D
holomorphic Chern--Simons on $X$~\cite{CostelloLi2016,CostelloLi2020}:
the BV action $S[A] = \int_X \mathrm{CS}_3(A) \wedge \Omega_X$ pairs
the $(0, \bullet)$-form part of a connection $A$ against the CY form,
producing a $(-1)$-shifted symplectic structure on the BV phase space.
The BV-quantised observables of 6D hCS form a holomorphic
$E_3$-factorisation algebra on $X$ which, by Costello--Li's
classification, is equivalent to the topological $E_3$-FA of Step~(b)
refined by $\bar\partial_X$ and the $\Omega_X$-induced CY pairing. The
holomorphic refinement is unconditional at $d = 1, 2$ and is the
chain-level CY-A$_d$ obstruction at $d = 3$, resolved
$(\infty, 1)$-categorically by Theorem~\ref{thm:cy-to-chiral-d3} and
chain-level at $d = 3$ by Theorem~\ref{thm:s3-framing-vanishes}
(three obstruction vanishings plus Bott periodicity
$\pi_3(B\Sp(2m)) = 0$).

The composite assembles
\begin{equation}
\label{eq:stage-one-output}
 \PhiFA_3(\cC) \;\in\; E_3\text{-}\HolFA(X)
\end{equation}
as an $E_3$-holomorphic factorisation algebra on $X$. Stage~1 is
canonical up to the contractible space of choices jointly controlled
by Kontsevich--Tamarkin $E_3$-formality (Step~(a), torsor under
$\GRT_1(\Q)$) and Costello--Gwilliam--Li holomorphic locality
(Step~(c), torsor under gauge transformations of the BV action); the
space of Stage~1 liftings is contractible as a fibration of two
contractible torsors over the initial datum $(\cC, X, \Omega_X)$. The
factorisation axiom
\begin{equation}
\label{eq:stage-one-factorisation}
 \PhiFA_3(\cC)(U \sqcup V) \;\simeq\;
 \PhiFA_3(\cC)(U) \otimes \PhiFA_3(\cC)(V)
\end{equation}
for disjoint opens $U, V \subset X$ is the local datum of the Stage-1
gluing; the Weiss-cover locality axiom (Definition~$6.1.2.1$ of
\cite{CostelloGwilliam} Vol.~I) supplies the global assembly.

\subsubsection{Stage~2: specialisation by factorisation homology over
$\Sigma_2$}
\label{subsec:fa-stage-two}

Stage~2 descends the $E_3$-holomorphic factorisation algebra of
\eqref{eq:stage-one-output} to an $E_1$-chiral algebra on a reference
curve $C \subset X$ by factorisation homology over a transverse
$2$-dimensional closed cycle $\Sigma_2 \subset X$ which fibres, in a
neighbourhood of $C$, as $\Sigma_2 \times C \hookrightarrow X$. The
tensor product of operads $E_n \otimes E_m$ (Boardman--Vogt tensor) is,
by Dunn's additivity theorem~\cite{Dunn1988} in its $\infty$-categorical
formulation due to Lurie~\cite{LurieHA}~Theorem~$5.1.2.2$,
\begin{equation}
\label{eq:dunn-lurie-e3-e2-e1}
 E_{n + m} \;\simeq\; E_n \otimes E_m,
 \qquad
 E_3 \;\simeq\; E_2 \otimes E_1.
\end{equation}
Equation~\eqref{eq:dunn-lurie-e3-e2-e1} identifies an $E_3$-algebra
with an $E_1$-algebra in $E_2$-algebras, or symmetrically an
$E_2$-algebra in $E_1$-algebras; the $E_2$-braided monoidal structure
on the inner category is the content of
Lurie~\cite{LurieHA}~\S$5.1.2.4$ (the characterisation of
$E_2$-algebras in $\Ch(k)$ as braided monoidal objects). Reading
\eqref{eq:dunn-lurie-e3-e2-e1} geometrically, the framed real
$6$-manifold $X$ near $C$ factors as a product $\Sigma_2 \times C$
(up to tubular neighbourhood), with $\Sigma_2$ supplying the transverse
$E_2$-directions and $C$ supplying the longitudinal $E_1$-direction;
factorisation homology over $\Sigma_2$
(\cite{LurieHA}~\S$5.5.3$; Francis~\cite{Francis2013} Theorem~$1.1$,
equivalently the push-forward formula \cite{Francis2013}~\S$3.4$)
integrates the $E_2$-directions against the transverse fibre and leaves
an $E_1$-factorisation algebra on $C$.

\begin{proposition}[Factorisation homology over a transverse fibre;
 Lurie~\cite{LurieHA}~\S$5.5.3$, Francis~\cite{Francis2013}
 Theorem~$1.1$ with push-forward \S$3.4$]
\label{prop:fa-transverse-integration}
Let $X$ be a smooth framed real $6$-manifold, $\Sigma_2 \subset X$ a
smooth closed framed $2$-submanifold, and $C \subset X$ a smooth framed
curve transverse to $\Sigma_2$ such that a tubular neighbourhood of
$\Sigma_2 \cap C$ inside $X$ is framed-diffeomorphic to
$\Sigma_2 \times C$. Then factorisation homology over $\Sigma_2$
defines an $(\infty, 1)$-functor
\begin{equation}
\label{eq:sigma2-integration}
 \int_{\Sigma_2} \colon E_3\text{-}\FA(X) \longrightarrow
 E_1\text{-}\FA(C),
\end{equation}
given by push-forward along the transverse projection
$\pi \colon \Sigma_2 \times C \to C$ of the $E_2$-factor of the
Dunn--Lurie splitting.
\end{proposition}

\begin{proof}[Sketch]
By Dunn--Lurie \eqref{eq:dunn-lurie-e3-e2-e1}, an
$E_3$-factorisation algebra on a framed real $6$-manifold is locally an
$E_2 \otimes E_1$-algebra. On the tubular neighbourhood $\Sigma_2
\times C$, the $E_2$-structure is carried by the $\Sigma_2$-direction
and the $E_1$-structure by the $C$-direction. Factorisation homology
of an $E_2$-algebra $A$ over a framed closed surface $\Sigma_2$ is a
chain complex $\int_{\Sigma_2} A \in \Ch(k)$
(Lurie~\cite{LurieHA}~\S$5.5.3$; Francis~\cite{Francis2013}~Theorem~$1.1$
for the general framed $n$-manifold case), functorial in $A$ as an
$E_2$-algebra. Applied fibrewise over $C$ to the $E_2$-factor of an
$E_2 \otimes E_1$-algebra, the $E_1$-structure on the $C$-factor is
preserved, yielding an $E_1$-factorisation algebra on $C$. The
explicit push-forward along $\pi$ is the product-manifold push-forward
of factorisation homology \cite{Francis2013}~\S$3.4$, equivalent on
the Ran-space presentation to the factorisation restriction of
\cite{FrancisGaitsgory2012} Theorem~$5.2.1$.
\end{proof}

The composite of Stage~1 with Stage~2 is the specialisation functor
\begin{equation}
\label{eq:spch-definition}
 \SpCh_{\Sigma_2, C} \colon E_3\text{-}\HolFA(X) \longrightarrow
 E_1\text{-}\HolFA(C),
 \qquad
 \cF \longmapsto \biggl(\int_{\Sigma_2} \cF\biggr) \Bigr|_{C},
\end{equation}
with the holomorphic refinement of \eqref{eq:sigma2-integration}
preserved by restriction to $C$. The two-stage factorisation of
$\Phi_3$ is then
\begin{equation}
\label{eq:phi3-two-stage}
 \Phi_3 \;=\; \SpCh_{\Sigma_2, C} \circ \PhiFA_3,
\end{equation}
reading the left-hand side as the CY-to-chiral functor at $d = 3$ and
the right-hand side as the composite of the Stage-1 factorisation-
algebra assembly and the Stage-2 specialisation.

\subsubsection{Each stage is itself a gluing}
\label{subsec:fa-each-stage-gluing}

The two-stage factorisation is a gluing in two layers. Stage~1 is
a factorisation-algebra assembly: the local data on Weiss-cover cells
of $X$ are $E_3$-algebras determined chart-wise by the BV quantisation
of 6D hCS, and the global assembly is the Weiss local-to-global axiom
applied to the disjoint-union closure of a disc atlas on $X$. Stage~2
is a specialisation assembly: the local data on Weiss-cover cells of
$C$ are $E_1$-algebras obtained by integrating the $E_3$-data of
Stage~1 against the transverse $\Sigma_2$-fibre, and the global
assembly is the Weiss local-to-global axiom on $C$ applied to a disc
atlas on $C$.

\begin{remark}[Six specialisations of a single Stage-$1$ output]
\label{rem:six-sigma2-C-specialisations}
A single CY$_3$ category $\cC$ produces, via Stage~1, a single
$E_3$-holomorphic factorisation algebra $\PhiFA_3(\cC)$ on a fixed
CY$_3$ target $X$. Different specialisation data
$(\Sigma_2, C), (\Sigma_2', C'), \dots$ on $X$ produce different
$E_1$-chiral shadows $\SpCh_{\Sigma_2, C}(\PhiFA_3(\cC))$,
$\SpCh_{\Sigma_2', C'}(\PhiFA_3(\cC))$, $\dots$; the multiplicity of
chiral images attached to $\cC$ is the multiplicity of Stage-2
specialisation cycles on $X$, not the multiplicity of $\Phi_3$-
applications to $\cC$. The six routes to $G(K3 \times E)$ are six
different $(\Sigma_2, C)$-specialisations of the single Stage-1 output
$\PhiFA_3(D^b(\Coh(K3 \times E)))$.
\end{remark}

\subsection{Holomorphic vs topological factors}
\label{sec:fa-hol-top}

\subsubsection{5D holomorphic Chern--Simons on $\R \times \C^2$}

The principal example of a factorisation algebra with distinct holomorphic
and topological factors is 5D holomorphic Chern--Simons on
$Y = \R_t \times \C^2_{z_1, z_2}$, introduced by
Costello~\cite{Costello2013} (arXiv:$1303.2632$), quantised to all
orders by Costello--Yagi~\cite{CostelloYagi2018}
(arXiv:$1810.01970$; two-author paper), with holomorphic-topological
factorisation controlled by Costello--Gwilliam Vol.~$2$~\S$4.10$
(arXiv:$2210.13036$)~\cite{CostelloGwilliam}. The manifold $Y$ is a
real $5$-manifold, product of the real $1$-dimensional topological
factor $\R_t$ and the complex $2$-dimensional (real $4$-dimensional)
holomorphic factor $\C^2_{z_1, z_2}$. The separation of holomorphic
from topological directions at the level of factorisation algebras is
Francis~\cite{Francis2013} (arXiv:$1303.0305$) Theorem~$1.1$:
$\Ch(k)$-valued factorisation algebras on a product of a topological
$p$-manifold and a holomorphic $q$-manifold are classified by
$E_p^{\mathrm{top}} \otimes E_q^{\mathrm{hol}}$-algebras.

Let $\fg$ be a finite-dimensional reductive Lie algebra with
non-degenerate invariant symmetric bilinear form $\tr$ (Killing form in
the simple case) and let $\Omega_{\C^2} = dz_1 \wedge dz_2$ be the
flat holomorphic volume form. The BV superfield is the total
$(0, \bullet) \otimes \bullet$-form
\begin{equation}
\label{eq:5d-hCS-field}
 \cA \;\in\; \Omega^{0, \bullet}(\C^2) \otimes \Omega^\bullet(\R_t)
 \otimes \fg[1],
 \qquad
 \cA \;=\; A_t(t, z, \bar z)\, dt
 \;+\; \sum_{i = 1, 2} A_{\bar z_i}(t, z, \bar z)\, d\bar z_i
 \;+\; \mathrm{ghosts} + \mathrm{antifields},
\end{equation}
including the ghost $c \in \Omega^{0, 0} \otimes \Omega^0 \otimes
\fg[1]$ and antifields in $\Omega^{0, \ge 1}(\C^2) \otimes
\Omega^{\ge 1}(\R_t) \otimes \fg$ as required by the BV shift. The BV
action of 5D hCS is
\begin{equation}
\label{eq:5d-hCS-action}
 S_{\mathrm{5D\,hCS}}[\cA] \;=\; \int_{\R_t \times \C^2}
 \Omega_{\C^2} \wedge \tr\!\left(
 \cA \wedge \bar\partial_{\C^2} \cA
 \;+\; \cA \wedge d_t \cA
 \;+\; \tfrac{2}{3}\, \cA \wedge \cA \wedge \cA
 \right),
\end{equation}
with $d_t$ the de Rham differential along $\R_t$ and
$\bar\partial_{\C^2}$ the Dolbeault operator on
$\C^2$~\cite{Costello2013,CostelloLi2020}. The classical equation of
motion is $(d_t + \bar\partial_{\C^2})\cA + \cA \wedge \cA = 0$: the
mixed flat-connection equation wedge-paired with $\Omega_{\C^2}$. The
normalisation $2/3$ in the cubic term arises from cyclic symmetrisation
of $\tr$ on $\fg^{\otimes 3}$, matching the BV-invariant master-equation
form of the Chern--Simons action.

\emph{Lorenz gauge.} The BV differential $Q_{\mathrm{BV}} = d_t +
\bar\partial_{\C^2} + [\cA, -]$ has a residual gauge freedom absorbed
by the mixed Lorenz condition
\begin{equation}
\label{eq:5d-hCS-lorenz-gauge}
 d_t^{\,\ast}\, A_t
 \;+\; \bar\partial_{\C^2}^{\,\ast}\!
 \bigl(A_{\bar z_1}\, d\bar z_1 + A_{\bar z_2}\, d\bar z_2\bigr)
 \;=\; 0,
\end{equation}
with $d_t^{\,\ast}, \bar\partial_{\C^2}^{\,\ast}$ the formal adjoints for
the flat Hermitian metric on $\C^2 \times \R_t$. On a contractible
open of $\R_t$ the admissible gauge slice sets $A_t = 0$, reducing the
topological factor to the $E_1$-multiplication of $\C^2$-observables
along the interval. On $\C^2$ the residual
$\bar\partial_{\C^2}$-harmonic sector carries the propagator
\begin{equation}
\label{eq:5d-hCS-propagator}
 \bigl\langle A_{\bar z_i}(z, t)\, A_{\bar z_j}(w, s)\bigr\rangle
 \;\propto\;
 \delta_{ij}\,\delta(t - s)\,\frac{1}{z_i - w_i},
\end{equation}
the pole source driving the all-orders Costello--Yagi
$\hbar$-expansion. The choice of Lorenz gauge fixes a BV-equivalent
representative of~\eqref{eq:5d-hCS-action}; factorisation-algebra
output is gauge-invariant.

The factorisation algebra of observables of 5D hCS admits a natural
tensor-product decomposition
\begin{equation}
\label{eq:5d-hCS-factorisation}
 \Obs^{\mathrm{5D\,hCS}}(\R_t \times \C^2) \;\simeq\;
 \Obs^{\mathrm{1D\,top}}(\R_t) \;\boxtimes\;
 \Obs^{\mathrm{4D\,hol}}(\C^2)
\end{equation}
with $\Obs^{\mathrm{1D\,top}}(\R_t)$ a locally constant $E_1$-factorisation
algebra on the topological factor $\R_t$ and $\Obs^{\mathrm{4D\,hol}}(\C^2)$
a \emph{holomorphic} $E_2^{\mathrm{hol}}$-factorisation algebra on
$\C^2$: the structure maps depend holomorphically on the centres of
$\C^2$-discs, with the two $\bar\partial_{\C^2}$-directions of
\eqref{eq:5d-hCS-field} supplying the two-dimensional OPE indexed by
holomorphic-disc configurations. The decomposition
\eqref{eq:5d-hCS-factorisation} is the factorisation-algebra
incarnation of Francis' holomorphic-topological splitting theorem
(Francis~\cite{Francis2013} Theorem~$1.1$, applied to the product of
a topological $\R_t$-factor and a holomorphic $\C^2$-factor).

At the level of the operad governing the $\Ch(k)$-valued algebra of
global observables, Kontsevich--Tamarkin $E_d$-formality on the
holomorphic factor identifies the holomorphic operad
$E_2^{\mathrm{hol}}(\C^2)$ with its topological shadow
$E_2^{\mathrm{top}}(\R^4)$ (Kontsevich~\cite{Kontsevich1999};
Tamarkin~\cite{Tamarkin2007}, 2D Deligne conjecture). Combined with
Dunn--Lurie~\cite{LurieHA}~Theorem~$5.1.2.2$ for the additivity of
topological $E_n$-operads, this produces the operad identification
\begin{equation}
\label{eq:dunn-lurie-e1-top-e2-hol}
 \mathrm{Op}\bigl(\Obs^{\mathrm{5D\,hCS}}\bigr)
 \;\simeq\;
 E_1^{\mathrm{top}} \otimes E_2^{\mathrm{hol}}
 \;\xrightarrow{\mathrm{formality}}\;
 E_1^{\mathrm{top}} \otimes E_2^{\mathrm{top}}
 \;\simeq\; E_3^{\mathrm{top}},
\end{equation}
with the arrow labelled ``formality'' encoding the
Kontsevich--Tamarkin equivalence $E_2^{\mathrm{hol}} \simeq
E_2^{\mathrm{top}}$ at the level of $\Ch(k)$-valued operads over
$\Q$. At the \emph{chain-level} sheaf geometry the holomorphic
structure of $\C^2$ is retained (equation~\eqref{eq:5d-hCS-factorisation}
is a strict external tensor product of sheaves on
$\R_t \times \C^2$); at the $(\infty, 1)$-categorical operad level
equation~\eqref{eq:dunn-lurie-e1-top-e2-hol} places global observables
in $E_3$-$\mathrm{Alg}(\Ch(k))$ after formality. The reading
\eqref{eq:dunn-lurie-e1-top-e2-hol} differs from the Stage-$2$
reading $E_3 \simeq E_2^{\Sigma_2} \otimes E_1^{C}$ of
\S\ref{subsec:fa-stage-two} only in which complex-geometric data
interprets the $E_2$-factor; the two readings coincide as abstract
$E_3$-algebras in $\Ch(k)$.

\subsubsection{Factorisation algebras on each factor separately}

On the topological factor $\R$, a locally constant factorisation algebra
is equivalent to an $E_1$-algebra in $\Ch(k)$ by
Lurie~\cite{LurieHA}~\S$5.5.4$ (specialisation to $n = 1$; equivalently,
the Costello--Gwilliam \cite{CostelloGwilliam} Vol.~I Theorem~$3.4.0.1$
identification of locally constant factorisation algebras on
$\R^n$ with $E_n$-algebras). The factor $\Obs^{\mathrm{1D\,top}}$
is the topological reduction of 5D hCS along $\R_t$: integrating the
BV action \eqref{eq:5d-hCS-action} against $dt$ and setting
$A_t = 0$ in Lorenz gauge leaves an effective $E_1$-algebra on $\R$
whose multiplication $\cA(U_1) \otimes \cA(U_2) \to \cA(U_1 \cup U_2)$
for disjoint $U_1, U_2 \subset \R_t$ is the $\hbar$-deformed
multiplication of the Yangian (Costello--Yagi $2018$
arXiv:$1810.01970$ Theorem~A, simply-laced bosonic case). Reading the Yangian as a topological
$E_1$-algebra on $\R_t$ alone discards the $\C^2$-OPE data that lives
on the holomorphic factor; the sharp joint statement identifies the
full factorisation algebra on $\R \times \C^2$ with the Yangian VOA
(Proposition~\ref{prop:nonabelian-5d-hcs-yangian-voa} below).

On the holomorphic factor $\C^2$, $\Obs^{\mathrm{4D\,hol}}$ is a
holomorphic $E_2$-factorisation algebra, equivalently a
$2$-dimensional vertex-algebra object in $\Ch(k)$ by the Francis
higher-dimensional analogue (arXiv:$1303.0305$ Theorem~$1.1$) of
Theorem~\ref{thm:fg-chiral-koszul}. The vertex-algebra output of
$\Obs^{\mathrm{4D\,hol}}$ is pinned by the $\bar\partial_{\C^2}$-OPE
of the $A$-field in \eqref{eq:5d-hCS-action}: the propagator
$\langle A_{\bar z_i}(z) A_{\bar z_j}(w)\rangle$ has a pole of order
$1$ on the diagonal $z = w$, and the residue carries the Yangian
coproduct structure in the Koszul-dual presentation supplied by
Proposition~\ref{prop:nonabelian-5d-hcs-yangian-voa}.

\begin{proposition}[Non-abelian 5D hCS $\to$ Yangian VOA, simply-laced
 bosonic; Costello--Yagi~\cite{CostelloYagi2018} (arXiv:$1810.01970$,
 two-author paper), Theorem~$4.1$]
\label{prop:nonabelian-5d-hcs-yangian-voa}
Let $\fg$ be a finite-dimensional simply-laced bosonic simple Lie
algebra (of type $A_n, D_n, E_6, E_7, E_8$) with non-degenerate
invariant Killing form. The BV-quantised factorisation algebra of
observables of 5D holomorphic Chern--Simons~\eqref{eq:5d-hCS-action}
on $\R_t \times \C^2$ with gauge algebra $\fg$, after compactification
along the topological factor $\R_t$, is equivalent to the Yangian VOA
$Y^{\mathrm{VOA}}(\fg)$ to all orders in $\hbar$:
\begin{equation}
\label{eq:5d-hcs-yangian-voa-allorders}
 \Obs^{\mathrm{5D\,hCS}, \mathrm{quant}}(\R_t \times \C^2, \fg)
 \Big|_{\R_t\text{-compact}}
 \;\simeq\;
 Y^{\mathrm{VOA}}(\fg).
\end{equation}
The perturbative expansion is convergent, not merely asymptotic, by
two rigidity inputs. First, Kontsevich--Tamarkin $E_2$-formality on the
holomorphic factor $\C^2$ (two-dimensional Deligne conjecture;
Kontsevich~\cite{Kontsevich1999}, Tamarkin~\cite{Tamarkin2007}) collapses
the $\hbar$-tower on $\Obs^{\mathrm{4D\,hol}}(\C^2)$ to its tree-level
OPE, with the Kontsevich weight formula bounding higher-loop Feynman
integrals. Second, in Lorenz gauge~\eqref{eq:5d-hCS-lorenz-gauge}
setting $A_t = 0$ on $\R_t$-opens, the only $1$D Feynman diagrams on
the topological factor are tadpoles, which cancel by the antisymmetry
of $\tr$ on $[\fg, \fg]$; this collapses the $\hbar$-tower on
$\Obs^{\mathrm{1D\,top}}(\R_t)$ to its tree-level $E_1$-multiplication.
The two collapses combined with the holomorphic-topological
splitting~\eqref{eq:5d-hCS-factorisation} determine the all-orders
factorisation algebra from its one-loop approximation. The one-loop
$r$-matrix of 5D hCS equals the Yangian $r$-matrix under the
identification Costello--Witten--Yamazaki~\cite{CostelloWittenYamazaki}
Theorem~$4.2$.

\emph{Cubic-versus-quadratic anomaly split.} The one-loop BV anomaly
on the $\tr(A\,\bar\partial A\,\bar\partial A\,\bar\partial A)$-wheel
splits as
\begin{equation}
\label{eq:5d-hCS-anomaly-split}
 A^{(1)}(\fg) \;=\; A^{(1)}_{\mathrm{cohom}} \;+\; A^{(1)}_{\mathrm{w.f.}},
 \qquad
 A^{(1)}_{\mathrm{cohom}}(\fg) \;=\; d^{abc}\, d^{abc},
 \qquad
 A^{(1)}_{\mathrm{w.f.}}(\fg) \;=\; -\,\frac{C_2(\fg)}{(2\pi)^3}
 \;=\; -\,\frac{2 h^\vee}{(2\pi)^3},
\end{equation}
with $d^{abc} = \tr(T^a \{T^b, T^c\})$ the totally symmetric rank-$3$
invariant tensor and $C_2(\fg) = 2 h^\vee$ the quadratic Casimir of
the adjoint (dual Coxeter number $h^\vee$). The cohomological
component $A^{(1)}_{\mathrm{cohom}}$ is scheme-independent and carries
the obstruction; the field-renormalisation component $A^{(1)}_{\mathrm{ren}}$
is scheme-dependent and absorbed into a BV-trivial counter-term. For
$\fg$ simply-laced of type $A_n, D_n, E_6, E_7, E_8$ the totally
symmetric invariant $3$-tensor vanishes on the adjoint (equivalently:
the only $\fg$-invariant in $S^3(\fg^\vee)$ is $0$ for simply-laced
Dynkin types; reducing to $\mathfrak{sl}_2$ gives $d^{abc} =
\frac{1}{4}\epsilon^{abc}$ antisymmetric, so $d^{abc} d^{abc} = 0$
upon symmetrised contraction). Consequently $A^{(1)}_{\mathrm{cohom}}
= 0$ on the simply-laced locus. Non-simply-laced bosonic $\fg$ of
types $B_n, C_n, F_4, G_2$ has $d^{abc} \neq 0$, obstructing the
$\hbar$-tower collapse beyond second order; the quantisation in that
regime is the twisted Yangian, and the all-orders identification
remains open.
\end{proposition}

\begin{remark}[Super gauge algebras: all-orders open]
\label{rem:5d-super-hcs-open}
The Costello--Yagi all-orders equivalence is established for
simply-laced \emph{bosonic} $\fg$. The super case
$\fg \in \{\widehat{\mathfrak{gl}}(m|n), \widehat{\mathfrak{osp}}(m|2n)\}$
-- the regime of the K$3$ Yangian envelopes $Y_\hbar(\mathfrak{so}(4, 20))$ and its Hodge-parity $\Z/2$-super-extension $Y_\hbar(\mathfrak{so}(4 \mid 20))$ (programme-specific, non-Kac; Kac $\mathfrak{osp}(4 \mid 20)$ excluded because its odd form is symplectic) and
the conifold two-chart super-Yangian; requires super
Kontsevich--Tamarkin formality, and
Ginzburg--Schedler~\cite{GinzburgSchedler2008} (Selecta Math.~$16$,
$2010$; arXiv:$0807.0174$) establish \emph{super Deligne} in the form
of super $E_2$-formality for Hochschild cochains of a Koszul Frobenius
super-algebra: the super Gerstenhaber bracket of degree $-1$ on super
Hochschild cochains lifts canonically to a super $E_2$-algebra
structure. The super $E_n$-formality at $n \ge 3$; in particular
the super $E_3$-formality required to lift the Hochschild cochain
bracket of a super CY$_3$ category to a super $E_3$-algebra, and the
super $E_4$-formality required for the $4$D holomorphic twist of super
hCS on $\C^2$ with super gauge $\fg$; is established only at
$E_2$-level and is open beyond.

The $1$-loop BV anomaly coefficient
$d^{abc}_{\mathrm{super}} = \mathrm{str}(T^a \{T^b, T^c\})$ vanishes
identically on $\mathfrak{gl}(1|1)$ by direct basis computation. With
bosonic Cartan pair $H, N$ and fermionic pair $\psi^+, \psi^-$
satisfying $\{\psi^+, \psi^-\} = N$, the non-zero super-symmetric
rank-$3$ contractions reduce to $\mathrm{str}(H \{N, N\}) +
\mathrm{str}(N \{H, N\}) + \mathrm{str}(\psi^+ \{\psi^-, H\}) +
\mathrm{str}(\psi^- \{\psi^+, H\})$, and each term vanishes: the
bosonic Cartan pair contributes $\mathrm{str}(H^2) = \mathrm{str}(N^2)
= 0$ by direct trace on the fundamental $(1|1)$-representation, and
the fermionic pair carries opposite supertrace signs $\mathrm{str}(\psi^+
\psi^-) = -\mathrm{str}(\psi^- \psi^+)$ which cancel under symmetric
contraction. The $1$-loop wheel on $\mathfrak{gl}(1|1)$ is therefore
verified. Higher-loop local cohomology
$H^1_{\mathrm{loc}}(\fg_{\mathrm{super}}, \cO_{\mathrm{loc}})^{\ge 2}$
is not automatically killed in the absence of super $E_n$-formality
for $n \ge 3$, and the all-orders super case remains open even on
$\mathfrak{gl}(1|1)$.
\end{remark}

The proposition is the sharpest instance of the
holomorphic-vs-topological factor decomposition: the topological factor
$\R$ supplies the $E_1$-multiplication of the Yangian, the holomorphic
factor $\C^2$ supplies the vertex-algebra/OPE structure, and the
Dunn--Lurie splitting \eqref{eq:5d-hCS-factorisation} packages the two
into a single $E_3$-factorisation algebra on $\R \times \C^2$.

\subsubsection{Gluing along stratified manifold structure}
\label{subsec:fa-stratified-gluing}

The product $\R_t \times \C^2$ is the simplest stratified manifold
with a holomorphic-topological decomposition: a single top stratum,
trivially stratified. General CY$_3$ targets carrying a 5D hCS
factorisation structure require the stratified-manifold machinery of
Costello--Gwilliam Vol.~$1$~\S$8$~\cite{CostelloGwilliam} (local-to-global
on locally conical stratified spaces), building on
Ayala--Francis--Tanaka~\cite{AyalaFrancisTanaka2017,AyalaFrancisTanaka2017Strat}
(\emph{Local structures on stratified spaces}, Adv.\ Math.~$307$;
\emph{Factorization homology of stratified spaces}, Selecta Math.~$23$),
which extends the factorisation-algebra local-to-global axiom to
stratified manifolds by replacing ordinary Weiss covers with
\emph{stratified Weiss covers} respecting the stratification.

\begin{definition}[Stratified Weiss cover;
 Costello--Gwilliam Vol.~$1$~\S$8$~\cite{CostelloGwilliam}
 Definition~$8.3.0.1$; Ayala--Francis--Tanaka
 \cite{AyalaFrancisTanaka2017} \S$2$--$3$, exit-path category of a
 conically smooth stratified space]
\label{def:stratified-weiss-cover}
Let $M$ be a conically smooth stratified manifold with stratification
$M = M_0 \sqcup M_1 \sqcup \cdots \sqcup M_r$, each $M_i$ a locally
closed smooth submanifold, the closures satisfying
$\overline{M_i} \subset \bigsqcup_{j \le i} M_j$. A collection $\fU$
of open subsets of $M$ is a \emph{stratified Weiss cover} if, for
every finite subset $S \subset M$, there exists $U_\alpha \in \fU$
with $U_\alpha \cap M_i \supset S \cap M_i$ simultaneously for every
stratum $M_i$, $i = 0, 1, \ldots, r$. Equivalently, $\fU$ is a Weiss
cover in the Grothendieck topology on the exit-path category of $M$
(\cite{AyalaFrancisTanaka2017Strat}~\S$3.4$).
\end{definition}

The extension of the factorisation local-to-global axiom to stratified
Weiss covers is Costello--Gwilliam Vol.~$1$~Theorem~$8.4.0.1$
\cite{CostelloGwilliam}, equivalently Ayala--Francis--Tanaka
\cite{AyalaFrancisTanaka2017Strat} Theorem~$1.1$ (factorisation
homology of stratified spaces is a homotopy invariant of the exit-path
category).
The effect for the present chapter is that the gluing of Stage~1 on a
stratified CY$_3$ target proceeds stratum by stratum, with the
top-dimensional stratum carrying the ambient 5D hCS factorisation
structure and finer (lower-dimensional) strata carrying corrections
supported on singular or exceptional loci.

\begin{example}[Local $\bP^2$ gluing via stratified Weiss covers]
\label{ex:local-p2-stratified-weiss}
Local $\bP^2 = \Tot(K_{\bP^2})$ is a complex $3$-fold (real $6$-manifold)
carrying the toric $T^2 = (\C^\times)^2$ action: the standard $T^2$
acts on $\bP^2$ by rescaling homogeneous coordinates modulo the
diagonal, and extends to the total space via the anti-canonical
character on the fibre of $K_{\bP^2}$. Stratify $\Tot(K_{\bP^2})$ by
locating each point either on the zero section $\bP^2$ (three toric
$T^2$-orbit types) or in its open complement:
\begin{equation}
\label{eq:local-p2-toric-strata}
 \Tot(K_{\bP^2}) \;=\;
 M_{\mathrm{pts}} \;\sqcup\; M_{\mathrm{edges}}
 \;\sqcup\; M_{\mathrm{bulk}} \;\sqcup\; M_{\mathrm{fibre}}.
\end{equation}
The four layers are:
\begin{enumerate}
\item[(i)] $M_{\mathrm{pts}}$ consists of the three $T^2$-fixed points
$p_1, p_2, p_3 \in \bP^2 \subset \Tot(K_{\bP^2})$, each an isolated
$0$-real-dimensional stratum on the zero section;
\item[(ii)] $M_{\mathrm{edges}}$ is the disjoint union of the three
toric edges $\bP^1_{ij} \setminus \{p_i, p_j\} \cong \C^\times$ on the
zero section (each the free orbit of a $1$-dimensional sub-$\C^\times
\subset T^2$), a $2$-real-dimensional locus;
\item[(iii)] $M_{\mathrm{bulk}} \cong (\C^\times)^2$ is the open
toric-bulk of the zero section: the single free $T^2$-orbit, the
$4$-real-dimensional complement of the three toric edges in
$\bP^2$;
\item[(iv)] $M_{\mathrm{fibre}} = \Tot(K_{\bP^2}) \setminus \bP^2$ is
the open complement of the zero section, a real $6$-manifold on which
$T^2 \times \C^\times_{\mathrm{fibre}}$ acts freely. The slice
$M_{\mathrm{fibre}}/\C^\times_{\mathrm{fibre}}$ is a $5$-real-dimensional
unit-link bundle over $\bP^2$ (topologically the $5$-sphere link of the
zero section), and we refer to $M_{\mathrm{fibre}}$ as the
$5$-real-dimensional generic fibre stratum in this link-quotient sense.
\end{enumerate}
The four layers partition the total space $\Tot(K_{\bP^2}) =
\bigsqcup_\bullet M_\bullet$ with closure relations
$\overline{M_{\mathrm{pts}}} = M_{\mathrm{pts}}$,
$\overline{M_{\mathrm{edges}}} = M_{\mathrm{pts}} \cup M_{\mathrm{edges}}$,
$\overline{M_{\mathrm{bulk}}} = \bP^2 = M_{\mathrm{pts}} \cup
M_{\mathrm{edges}} \cup M_{\mathrm{bulk}}$, and $\overline{M_{\mathrm{fibre}}}
= \Tot(K_{\bP^2})$.

A stratified Weiss cover $\fU^{\mathrm{strat}}$ of local $\bP^2$
assigns, to each finite configuration $S \subset \Tot(K_{\bP^2})$, an
open $U \in \fU^{\mathrm{strat}}$ whose intersection with each of the
four strata in \eqref{eq:local-p2-toric-strata} contains the
corresponding slice $S \cap M_\bullet$. The stratum-by-stratum
factorisation algebra reads: $M_{\mathrm{pts}}$ contributes the
McKay-quiver CoHA centred on each toric fixed point's local $\C^3$
chart via the Bridgeland--King--Reid equivalence $D^b(\Coh([\C^3/\Z_3]))
\simeq D^b(\mathrm{mod}\text{-}\mathrm{Beil}_3)$ for the Beilinson
quiver $\mathrm{Beil}_3$ of $\bP^2$ (\S\ref{sec:gluing:toric-cech},
\S$3.2$, Szendr\H{o}i--Nakajima); $M_{\mathrm{edges}}$ contributes
three $\bP^1$-normal-bundle chiral algebras, each an
$E_2$-factorisation algebra on an edge curve (Costello--Li
\cite{CostelloLi2016} $4$D hCS reduction over the edge);
$M_{\mathrm{bulk}}$ and $M_{\mathrm{fibre}}$ together recover the
bulk 5D hCS reduction of local $\bP^2$ at $T^2$-generic parameter,
matching the BKR equivalence on the zero section (Beilinson quiver of
$\bP^2$). The toric stratification is the reason the CoHA of local
$\bP^2$ decomposes, via the Weiss-cover gluing of
Definition~\ref{def:stratified-weiss-cover}, into the three
quiver-CoHAs at $M_{\mathrm{pts}}$ plus the three edge corrections at
$M_{\mathrm{edges}}$ plus the bulk contribution at $M_{\mathrm{bulk}}
\sqcup M_{\mathrm{fibre}}$, and not merely into the open-versus-zero-
section pair $\Tot(K_{\bP^2}) \setminus \bP^2 \sqcup \bP^2$ which
ignores the toric equivariance structure.
\end{example}

\begin{remark}[Connection to Costello--Paquette $\C^3$ M-theory]
\label{rem:costello-paquette-c3-m-theory}
The Costello--Paquette programme (Costello--Paquette arXiv:$1810.06490$
for $\C^3$ and arXiv:$2009.04834$ for the Omega-background extension)
identifies the twisted M-theory factorisation algebra on $\R \times
\C^2 \times S^1$ with the $\R$-quantised 5D hCS theory of
Proposition~\ref{prop:nonabelian-5d-hcs-yangian-voa}, specialised to
$\fg = \widehat{\fgl}_\infty$ and twisted by an Omega-background on
the $\C^2$ holomorphic factor. The stratified Weiss cover of the
M-theory $\C^3$-geometry decomposes as the toric $T^3$-stratification
of $\C^3$ (one open stratum plus three $\C^2 \subset \C^3$ toric-
divisor strata plus three $\C \subset \C^2$ toric-edge strata plus the
origin $0 \in \C^3$), and the reduction to 5D hCS on $\R \times \C^2$
amounts to compactifying the $S^1$-factor and restricting to the
generic $T^2 \subset T^3$-invariant slice. The $\widehat{\fgl}_\infty$
Yangian VOA arising in this reduction is the M-theoretic extension of
the finite-rank Yangian VOA of
Proposition~\ref{prop:nonabelian-5d-hcs-yangian-voa}, with the
$T^3$-stratified Weiss-cover gluing supplying the degeneration
structure across the codimension-$1$ toric-divisor walls.
\end{remark}

\begin{remark}[Strict super-factorisation across strata]
\label{rem:strict-super-factorisation-stratified}
The strict super-factorisation axiom of
Remark~\ref{rem:strict-super-factorisation} extends to stratified
Weiss covers: the Koszul super-braiding in the structure maps
\eqref{eq:prefactorisation-structure-map} restricted to each stratum
remains a strict equality in $\Ch(\sVect)$, and the
stratum-$\Z/2$-grading compatibility along codimension-$1$ strata is a
strict equality of interchange isomorphisms rather than a coherence
homotopy. The stratification does not promote the super-factorisation
axiom to a coherence datum. The reason the K3 $\times$ E fibre
stratification (compact fibre $K3 \times \{\mathrm{pt}\}$ over an
elliptic base $\{\mathrm{pt}\} \times E$) carries a strict
$\Z/2$-grading on the global chiral algebra is that each of the two
strata is itself a factorisation algebra in $\Ch(\sVect)$ with the
strict interchange, and the gluing along the stratum boundary uses
only the strict $1$-categorical symmetry of $\Ch(\sVect)$; no
$(\infty, 1)$-categorical homotopy datum is required at this
step. The same strict extension applies to the local $\bP^2$ toric
stratification of Example~\ref{ex:local-p2-stratified-weiss} with
the three fixed-point $\C^3$-CoHAs carrying strict $\Z/2$-gradings
inherited from the conifold prototype of \S\ref{sec:gluing:toric-cech}, \S$2.1$.
\end{remark}

\subsection{Section summary}

Sections~\ref{sec:fa-weiss-covers}--\ref{sec:fa-hol-top} establish the
factorisation-algebra gluing mode of \S\ref{sec:gluing:factorisation} at four elementary levels.
The Weiss-cover criterion (Definition~\ref{gl:def:weiss-cover}) replaces
ordinary open covers with finite-point-configuration covers; a
two-chart atlas is not a Weiss cover
(Proposition~\ref{prop:two-chart-not-weiss}), and the disjoint-union
closure of a subordinate polydisc basis is the Weiss refinement
\eqref{eq:weiss-refinement}. The Ran space of a curve indexes chiral
observables as a factorisation cosheaf, and the Francis--Gaitsgory
chiral Koszul duality (Theorem~\ref{thm:fg-chiral-koszul}) identifies
chiral algebras on a curve with factorisation cosheaves on the Ran
space. Proposition~\ref{prop:finite-dwr-ran-rees-ordered-hall-layer}
states the finite DWR/Ran hCS--Rees-Hall comparison on this same
finite-configuration site: finite descent is over an ordered $E_1$ Hall
layer, the Hall product is the Rees Thom--Sebastiani product, and the
$E_2$ centre or Drinfeld double requires its own pairing/completion
operation. The higher-dimensional Francis $E_n$-analogue
(Theorem~\ref{thm:francis-en-koszul}) underlies the CY$_3$ construction.
The two-stage factorisation $\Phi_3 = \SpCh_{\Sigma_2, C} \circ
\PhiFA_3$ (equations~\eqref{eq:stage-one-output}--\eqref{eq:phi3-two-stage})
assembles Stage~1 by the Costello--Gwilliam factorisation-envelope
refinement of Kontsevich--Tamarkin $E_3$-formality and descends to
Stage~2 by factorisation homology over a transverse $\Sigma_2$-fibre
(Proposition~\ref{prop:fa-transverse-integration}); each stage is
itself a Weiss-cover gluing. The 5D holomorphic Chern--Simons BV action $S_{\mathrm{5D\,hCS}}[\cA]
= \int_{\R_t \times \C^2} \Omega_{\C^2} \wedge \tr(\cA \wedge
\bar\partial_{\C^2} \cA + \cA \wedge d_t \cA + \tfrac{2}{3}\, \cA
\wedge \cA \wedge \cA)$ of~\eqref{eq:5d-hCS-action}, in mixed Lorenz
gauge~\eqref{eq:5d-hCS-lorenz-gauge}, produces a factorisation algebra
on $\R_t \times \C^2$ decomposing~\eqref{eq:5d-hCS-factorisation}
into a topological $E_1$-factor on $\R_t$ and a holomorphic
$E_2^{\mathrm{hol}}$-factor on $\C^2$ by Francis's
holomorphic-topological splitting~\cite{Francis2013} Theorem~$1.1$;
Kontsevich--Tamarkin formality then identifies the operad with $E_3$
via~\eqref{eq:dunn-lurie-e1-top-e2-hol}. Non-abelian 5D hCS converges
to the Yangian VOA at all orders for simply-laced bosonic $\fg$ of
types $A_n, D_n, E_6, E_7, E_8$ by the cubic anomaly $d^{abc} d^{abc}
= 0$ (the totally symmetric invariant $3$-tensor vanishes on
simply-laced adjoint) and the scheme-dependent $2 h^\vee$-counter-term
(Proposition~\ref{prop:nonabelian-5d-hcs-yangian-voa}), and the super
case remains open except for the verified $1$-loop cancellation on
$\mathfrak{gl}(1|1)$ (Remark~\ref{rem:5d-super-hcs-open}). Stratified Weiss covers
(Definition~\ref{def:stratified-weiss-cover}; Ayala--Francis--Tanaka
\cite{AyalaFrancisTanaka2017,AyalaFrancisTanaka2017Strat}) extend the
local-to-global axiom to stratified targets; the local $\bP^2$ toric
$T^2$-stratification (Example~\ref{ex:local-p2-stratified-weiss}) has
four layers $M_{\mathrm{pts}} \sqcup M_{\mathrm{edges}} \sqcup
M_{\mathrm{bulk}} \sqcup M_{\mathrm{fibre}}$ of real dimensions
$0, 2, 4, 5$ respectively (three isolated fixed points of dim $0$;
three toric $\bP^1$-edges minus fixed points of dim $2$; one open
toric bulk $(\C^\times)^2 \subset \bP^2$ of dim $4$; and the
fibre-direction-open link of dim $5$), each carrying a stratum-specific
factorisation algebra whose gluing reproduces the CoHA of local
$\bP^2$. The
Costello--Paquette $\C^3$ M-theory picture (arXiv:$1810.06490$,
arXiv:$2009.04834$; Remark~\ref{rem:costello-paquette-c3-m-theory})
identifies the $T^3$-stratified $\C^3$-geometry of twisted M-theory with
the $S^1$-compactification of 5D hCS on $\R \times \C^2$, specialised
to $\fg = \widehat{\fgl}_\infty$. The strict super-factorisation
axiom in $\Ch(\sVect)$
(Remarks~\ref{rem:strict-super-factorisation} and
\ref{rem:strict-super-factorisation-stratified}) is $1$-categorical
and is inherited strictly by stratified Weiss covers: the
$\Z/2$-grading on the super-Yangian arising from the conifold two-chart
gluing (\S\ref{sec:gluing:toric-cech}, \S$2.1$), on the global chiral algebra of
K3 $\times$ E, and on the three fixed-point CoHAs of local $\bP^2$ is
a strict equality, not a coherence datum. \S\ref{sec:gluing:wallcrossing} turns to
wall-crossing gluing, stratum-independent and complementary to
Weiss-cover locality.
